\section{Structural Isomorphism with HPA: Protocol-Layer Mapping and Three-Channel Organization}

\subsection{Alignment of Protocol-Layer Objects}

Aligning the three-layer framework in this paper with the HPA protocol layer yields the following mapping:
\begin{itemize}
  \item \textbf{Scan operator:} $T_g$ (compact-group translation) corresponds to HPA's unit scan operator (with iteration count as operational timescale);
  \item \textbf{Readout kernel:} window $\ind_W$ and $\pi_\varepsilon$ correspond to HPA projective readout (an abstraction of finite-resolution POVM/kernel readout);
  \item \textbf{Canonical coding:} Sturmian/Ostrowski/Zeckendorf readout corresponds to HPA canonical discrete coordinates (converting phase orbits into auditable strings);
  \item \textbf{Renormalization:} substitution and blocking operators correspond to HPA cross-scale stability mechanisms (fixed points, attractors, and frequency closure).
\end{itemize}

\subsection{Three-Channel Semantics of \texorpdfstring{$(\varphi,\pi,e)$}{(varphi, pi, e)}}

Within one protocol family, $\varphi$, $\pi$, and $e$ correspond to three mutually commutable structural constraints:
\begin{itemize}
  \item \textbf{$\varphi$ channel:} anti-locking in the golden branch (approximation extremality), together with Zeckendorf admissibility/golden-mean subshift grammar, organizes the stable type space;
  \item \textbf{$\pi$ channel:} closure/periodic structure of phase space and cyclic admissibility (endpoint gluing) induce additional refinements and boundary/cycle splitting;
  \item \textbf{$e$ channel:} Abel/Euler-type normalization and exponential completion provide analytic stability and auditable conditions such as ``no interior poles in the unit disk.''
\end{itemize}

In this sense, the present paper gives a complete mathematical instance showing that HPA ``multichannel stability'' can be placed in standard ergodic-symbolic-analytic normalization language.
